% Source: Supplement: implication graph and SCC; https://cp-algorithms.com/graph/2SAT.html
% Review and validation: docs/validation.md.
\subsection{2-SAT：二元布尔约束}
\label{algo:two-sat}
\AlgorithmInfo
    {判断若干个二文字析取子句能否同时满足，并构造一组布尔赋值。}
    {有 $n$ 个变量、$m$ 个子句，时间、空间 $O(n+m)$。}
    {变量从 1 开始；必须使用本章 SCC 的逆拓扑编号约定。}

\subsubsection*{原理}
每个变量拆成两个文字 $(x_i=0)$、$(x_i=1)$。子句 $A\lor B$ 等价于两条蕴含 $\neg A\Rightarrow B$、$\neg B\Rightarrow A$，据此建立有向图。

若某变量的两个文字在同一个 SCC，选择其中任意一个都会推出另一个，必然矛盾。反之，当所有变量都不满足这一冲突条件时，可以按缩点图的拓扑关系构造解。由于本章 Tarjan 的 SCC 编号为逆拓扑序，取
\[
    x_i=1\quad\Longleftrightarrow\quad
    bel_{(x_i=1)}<bel_{(x_i=0)}.
\]
若更换成顺拓扑编号的 SCC 实现，比较方向也必须随之调整。

\subsubsection*{常用建模}
强制 $A$ 为真可加入 $A\lor A$；$A\Rightarrow B$ 对应 $\neg A\lor B$，需要同时加入其逆否命题的边。$A=B$ 加入 $A\Rightarrow B$ 与 $B\Rightarrow A$；$A\ne B$ 加入 $A\lor B$ 和 $\neg A\lor\neg B$。三个文字任意组成的析取不能直接当作普通 2-SAT 子句。

\subsubsection*{使用教程}
先复制 SCC 模板。\InlineCode{either(x,a,y,b)} 表示 $(x_x=a)\lor(x_y=b)$；\InlineCode{imply(x,a,y,b)} 表示前者成立则后者必须成立，内部会补齐逆否边；\InlineCode{force(x,a)} 强制值。\InlineCode{solve()} 返回真后读 \InlineCode{ans[i]}，为假则无解。
\begin{lstlisting}[style=mtycpp,numbers=none]
TwoSAT sat;
sat.init(3);
sat.either(1,true,2,false);
sat.imply(2,true,3,true);
sat.force(1,true);
bool ok=sat.solve();
if(ok){int x=sat.ans[3];}
\end{lstlisting}

\needspace{7\baselineskip}
\subsubsection*{模板代码}
\CodeFile{code/graphs/two_sat.cpp}
